\section{Investment Portfolio}

The I2.0 investment portfolio has seven components. Each component is analyzed independently for capital cost, annual benefit, NPV at 7\% discount, and benefit-cost ratio. The portfolio total is then phased across a 30-year implementation horizon. Table~\ref{tab:portfolio} summarizes the complete portfolio.

\subsection{Component 1: Managed Freight Lanes (E.1)}

The managed freight lane program covers seven T1 corridors: I-95 (NE--FL), I-80 (NY--CA), I-10 (FL--CA), I-40 excluded (see Section~2), I-70 (east segment, PA--CO), I-35 (MN--TX), I-5 (CA--WA), and I-90 (Boston--Seattle). The design specification calls for two dedicated lanes per direction, separated from general traffic by flexible delineators, with access restricted to commercial vehicles and managed through dynamic tolling \citep{ROUTE_E1,TRB_ML2012}.

\begin{itemize}
  \item \textbf{Capital cost}: \$121 billion over 15 years (Phase 1 and 2)
  \item \textbf{Annual benefit}: \$12.7 billion per year (freight time savings \$8.2B, shipper reliability premium \$4.5B \citep{ROUTE_C1,ROUTE_E1})
  \item \textbf{NPV at 7\%}: \$86.4 billion
  \item \textbf{B/C ratio}: 2.3:1
\end{itemize}

The \$12.7 billion annual benefit is derived from two streams. The freight time savings stream (\$8.2B/yr) uses the FHWA Freight Performance Measures baseline PTI of 1.86 on I-80 and 1.74 on I-35, reduced to 1.15 on managed lane segments, applied to the annual commercial vehicle volumes and a standard value-of-time estimate of \$42.14 per hour (FHWA 2023 guidance) \citep{FHWA_freight2023}. The shipper reliability premium (\$4.5B/yr) is derived from the Track C analysis of shipper buffer stocks maintained against PTI variance \citep{ROUTE_C1}: the managed lane PTI of 1.15 allows shippers to reduce transit time buffers from 80 hours to 48 hours, releasing inventory capital valued at current commercial lending rates.

\subsection{Component 2: Diamond Interchange Upgrades (B.4)}

Track B identified 9 of 15 T1/T1 intersections as k=1 nodes: they have no alternate route in the event of the primary interchange being closed \citep{ROUTE_B4}. The remaining 6 T1/T1 intersections have partial redundancy (k=2) but at-grade mixing that creates bottleneck conditions under freight demand. All 15 receive diamond interchange upgrades.

\begin{itemize}
  \item \textbf{Capital cost}: \$4.5 billion (average \$300M per interchange; range \$180M--\$520M depending on existing footprint and urban land costs)
  \item \textbf{Annual benefit}: \$3.2 billion per year (incident cost reduction \$1.8B, travel time reliability \$1.4B \citep{ROUTE_B4})
  \item \textbf{NPV at 7\%}: \$41.3 billion
  \item \textbf{B/C ratio}: 2.76:1 --- the highest per-dollar B/C in the portfolio, reflecting the severity of the k=1 condition
\end{itemize}

\subsection{Component 3: Compound Exposure Hardening (B.3)}

Track B identified 11 compound exposure corridors --- those scoring above threshold on both the climate exposure dimension (D1 $>$ 6) and the structural fragility dimension (B1 $>$ 7) \citep{ROUTE_B3}. The top 7, ranked by combined risk score, receive hardening investment in Phase 1. Gulf Coast I-10 is the highest-priority; its D1 score of 9.1 (projected to 2050) reflects combined sea-level rise, hurricane wind loading, and storm surge flood risk \citep{ROUTE_D1}.

\begin{itemize}
  \item \textbf{Capital cost}: \$18 billion (range \$1.2B to \$4.8B per corridor depending on length and hazard type)
  \item \textbf{Annual benefit}: \$3.9 billion per year in closure cost reduction \citep{ROUTE_D2}
  \item \textbf{NPV at 7\%}: \$37.8 billion
  \item \textbf{B/C ratio}: 3.1:1 --- highest per-dollar B/C because hardening addresses two risk dimensions simultaneously
\end{itemize}

\subsection{Component 4: Missing Link Construction (B.1)}

The five proposed corridors meeting the T1 rubric threshold represent \$88 billion in new construction or full-standard upgrade investment. Their individual cases are:

\begin{itemize}
  \item \textbf{I-69 completion}: \$22 billion for approximately 400 miles of new construction (segment from Texarkana TX to Indianapolis IN not yet completed to interstate standard). B/C ratio at 7\% discount: 1.4:1. At social discount rate of 5\% (standard for large infrastructure with multi-generational benefits): B/C 2.1:1 \citep{ROUTE_B1,ROUTE_C2}.
  \item \textbf{I-31/US-287 upgrade}: \$18 billion; upgrade 920 miles of US-287 to interstate standard (Amarillo TX to Havre MT). B/C 2.0:1, with a military access premium reflected in the B4 dimension score \citep{ROUTE_B1}.
  \item \textbf{I-29S/US-83 upgrade}: \$16 billion; 840 miles of US-83 from Laredo TX to the existing I-29 terminus at Sioux Falls SD. B/C 1.8:1; agricultural economy connectivity (C3) and tornado corridor resilience (D1) are primary drivers \citep{ROUTE_B1}.
  \item \textbf{I-92/US-2 upgrade}: \$20 billion; 1,050 miles of US-2 from Seattle WA to Minot ND. B/C 1.4:1; rural connectivity and northern tier redundancy are primary drivers \citep{ROUTE_B1}.
  \item \textbf{I-70W/US-50}: \$12 billion; 560 miles of US-50 from Reno NV to Pueblo CO upgraded to interstate standard. B/C 1.2:1 at 7\% discount; resilience value (Donner alternative) is the primary justification. Without the resilience premium, the corridor does not meet the investment threshold on demand alone \citep{ROUTE_D2}.
\end{itemize}

\subsection{Component 5: Intermodal Integration}

Intermodal terminal investments on T1 and T2 corridors --- expanded staging capacity, crane upgrades, and terminal access road improvements at 24 existing intermodal facilities --- address the last-mile friction that erodes the managed lane PTI gains.

\begin{itemize}
  \item \textbf{Capital cost}: \$12 billion
  \item \textbf{Annual benefit}: \$2.1 billion per year (dwell time reduction, empty repositioning cost reduction \citep{FHWA_freight2023})
  \item \textbf{NPV at 7\%}: \$17.5 billion
  \item \textbf{B/C ratio}: 1.8:1
\end{itemize}

\subsection{Component 6: EV Charging Infrastructure}

DC Fast Charging deployment at T1 rest areas, at 50-mile maximum spacing, supports the transition of Class 8 freight to battery-electric and fuel cell powertrains. At current commercial EV truck adoption rates (approximately 4\% of new Class 8 registrations in 2025), the return is modest; the benefit grows as adoption increases through 2035.

\begin{itemize}
  \item \textbf{Capital cost}: \$8 billion (approximately 3,200 charging stations at \$2.5M average installed cost, including grid interconnection upgrades)
  \item \textbf{Annual benefit}: \$0.9 billion per year at 2025 EV adoption rates; projected \$2.8B/yr at 2035 adoption rates \citep{FHWA_NHS2023}
  \item \textbf{NPV at 7\%}: \$9.2 billion (using 2025--2055 benefit trajectory)
  \item \textbf{B/C ratio}: 1.2:1 at current adoption; 2.1:1 at 2035 adoption
\end{itemize}

\subsection{Component 7: Transit Integration Layer (F.1)}

The transit integration layer --- bus platforms, passenger parking, and shelters at 9 T1/T1 diamond hubs and approximately 50 T1/T2 regional stops --- is fully described in the companion F.1 paper. Its NPV is quantified in F.2.

\begin{itemize}
  \item \textbf{Capital cost}: \$2 billion
  \item \textbf{Annual benefit}: quantified in F.2 (intercity bus corridor analysis)
\end{itemize}

\subsection{Portfolio Summary and Phasing}

\begin{table}[h]
\centering
\caption{I2.0 Investment Portfolio Summary}
\label{tab:portfolio}
\begin{tabular}{lrrrr}
\toprule
Component & Capital (\$B) & Annual Benefit (\$B/yr) & NPV at 7\% (\$B) & B/C \\
\midrule
1. Managed freight lanes & 121.0 & 12.7 & 86.4 & 2.3:1 \\
2. Diamond interchange upgrades & 4.5 & 3.2 & 41.3 & 2.76:1 \\
3. Compound exposure hardening & 18.0 & 3.9 & 37.8 & 3.1:1 \\
4a. I-69 completion & 22.0 & 2.1 & 8.5 & 1.4:1 \\
4b. I-31/US-287 upgrade & 18.0 & 2.4 & 13.6 & 2.0:1 \\
4c. I-29S/US-83 upgrade & 16.0 & 2.0 & 11.0 & 1.8:1 \\
4d. I-92/US-2 upgrade & 20.0 & 1.9 & 9.5 & 1.4:1 \\
4e. I-70W/US-50 & 12.0 & 0.9 & 5.6 & 1.2:1 \\
5. Intermodal integration & 12.0 & 2.1 & 17.5 & 1.8:1 \\
6. EV charging infrastructure & 8.0 & 0.9 & 9.2 & 1.2:1 \\
7. Transit integration (F.1) & 2.0 & (F.2) & (F.2) & (F.2) \\
\midrule
\textbf{Total (ex-transit)} & \textbf{251.5} & \textbf{32.1} & \textbf{240.4} & \textbf{2.2:1} \\
\textbf{Total (incl. transit)} & \textbf{253.5} & --- & --- & --- \\
\bottomrule
\end{tabular}
\end{table}

\noindent Note: Component B/C ratios are calculated independently; the aggregate B/C reflects positive complementarity among components.

\textbf{Implementation phasing} distributes the \$253 billion over three decades:
\begin{itemize}
  \item \textbf{Phase 1 (2025--2035), \$89 billion}: Managed lanes on the three highest-B/C corridors (I-95, I-35, I-5); all 15 diamond interchange upgrades; all 7 compound exposure hardening programs; transit integration layer. This phase delivers the highest-return investments first and hardens the most vulnerable nodes before the 2035 climate risk inflection.
  \item \textbf{Phase 2 (2035--2045), \$96 billion}: I-69 completion; I-31/US-287 and I-29S/US-83 upgrades; managed lanes on remaining T1 corridors (I-80, I-10, I-70, I-90). Missing link construction requires the longest lead time for right-of-way acquisition and environmental review.
  \item \textbf{Phase 3 (2045--2055), \$68 billion}: I-92/US-2 upgrade; I-70W/US-50; intermodal integration; EV charging completion. Third-phase investments are less time-sensitive because their benefits are either tied to EV adoption trajectories or serve corridors with lower current demand.
\end{itemize}
